% =============================================================
% บทที่ 1 - บทนำ
% โครงงาน: การออกแบบและพัฒนาเกมเดินเรือด้วยระบบสั่งการด้วยเสียงเพื่อส่งเสริมทักษะการสื่อสาร
%          (Design and Development of a Voice-Controlled Ship Navigation Game to Enhance Communication Skills)
% =============================================================

% -----------------------------------------------------------
\section{ความเป็นมาและความสำคัญของปัญหา}
% -----------------------------------------------------------
ในปัจจุบัน ภาษาอังกฤษมีบทบาทสำคัญอย่างยิ่งในการเป็นเครื่องมือสื่อสารระดับสากล ทว่าการเรียนการสอนภาษาอังกฤษในห้องเรียนโดยทั่วไปยังคงเผชิญกับข้อจำกัดในการสร้างสภาพแวดล้อมที่เอื้อต่อการนำไปใช้งานจริง\cite{sadigzade2025} ผู้เรียนจำนวนมากประสบปัญหาความวิตกกังวล (Language Anxiety) และกลัวความผิดพลาดเมื่อต้องเปล่งเสียงเพื่อสื่อสาร ส่งผลให้ขาดความมั่นใจและไม่สามารถเชื่อมโยงทักษะทางภาษาเข้ากับสถานการณ์ในชีวิตจริงหรือการโต้ตอบแบบทันท่วงทีได้\cite{varma2025} 
ที่ผ่านมา มีความพยายามนำสื่อดิจิทัลและเกมมาประยุกต์ใช้เพื่อการเรียนรู้ภาษา (Game-Based Language Learning) แต่ยังคงพบกับความท้าทายสำคัญ\cite{sadigzade2025, sanchez2010, allain2015} งานวิจัยชี้ให้เห็นว่าแอปพลิเคชันหรือเกมเพื่อการเรียนรู้ภาษาแบบดั้งเดิมมักมีความหนาแน่นทางไวยากรณ์และการท่องจำคำศัพท์ ทำให้รูปแบบการเล่นจำเจอยู่กับการคลิกเลือกคำตอบหรือการอ่านข้อความ ในขณะที่เกมที่เน้นความสนุกสนานก็มักไม่ได้บังคับให้ผู้เล่นต้องเปล่งเสียงพูดออกมาจริง ๆ\cite{varma2025} ช่องว่างนี้แสดงให้เห็นถึงโอกาสในการสร้างสรรค์เกมที่สามารถบูรณาการการฝึกฝนทักษะการพูดเข้ากับบริบทที่น่าตื่นเต้นและท้าทาย เพื่อกระตุ้นให้เกิดการสื่อสารอย่างเป็นธรรมชาติ\cite{zargham2024}
โครงงานการออกแบบและพัฒนาเกมเดินเรือด้วยระบบสั่งการด้วยเสียงเพื่อส่งเสริมทักษะการสื่อสารจึงถูกนำเสนอขึ้น เพื่อพัฒนาแนวทางการเรียนรู้ภาษาอังกฤษผ่านสภาพแวดล้อมเสมือนจริง (Immersive Environment) ผู้เล่นจะรับบทเป็น "กัปตันเรือ" ที่ต้องเผชิญกับสถานการณ์จำลองการเดินเรือฝ่าสภาพแวดล้อมที่ท้าทาย โดยภารกิจของผู้เล่นไม่ใช่เพียงการควบคุมทิศทางผ่านหน้าจอ แต่คือการใช้ "เสียง" เป็นเครื่องมือหลักในการควบคุมและสั่งการลูกเรือปัญญาประดิษฐ์ (AI Crew) ในตำแหน่งต่าง ๆ เช่น พลต้นหน (Navigator) และพลถือท้ายเรือ (Helmsman) เพื่อบังคับเรือไม่ให้ชนสิ่งกีดขวางและนำเรือไปสู่เป้าหมายอย่างปลอดภัย

% -----------------------------------------------------------
\section{นิยามศัพท์เฉพาะ}
% -----------------------------------------------------------

\begin{enumerate}[label=\arabic*)]
    \item \textbf{เกมเพื่อการเรียนรู้ (Serious Game)} หมายถึง เกมแนวจำลองสถานการณ์ที่ถูกออกแบบมาเพื่อใช้เป็นสื่อการเรียนรู้แบบดื่มด่ำ (Immersive Learning) ที่มีความสมจริง
    \item \textbf{การรู้จำเสียง (Speech-to-Text / Speech Recognition)} หมายถึง เทคโนโลยีการแปลงคำสั่งเสียงพูดภาษาอังกฤษของผู้เล่นให้เป็นข้อความ เพื่อนำไปตรวจความถูกต้องและประมวลผลเป็นคำสั่งควบคุมเรือ
    \item \textbf{ตัวกรองทางอารมณ์ (Affective Filter)} หมายถึง อุปสรรคทางอารมณ์ เช่น ความกลัว ความวิตกกังวล และความขาดความมั่นใจ ซึ่งขัดขวางการเรียนรู้และการพูดภาษาที่สอง
    \item \textbf{ลูกเรือปัญญาประดิษฐ์ (AI Crew)} หมายถึง ตัวละคร NPC (Non-Player Character) ในเกมที่รับคำสั่งเสียงจากผู้เล่นและให้เสียงตอบกลับ (Voice Feedback) เช่น การทวนคำสั่งหรือขอให้ย้ำคำสั่ง
    \item \textbf{สภาพแวดล้อมเสมือนจริง (Immersive Environment)} หมายถึง สภาพแวดล้อมจำลองภายในเกมที่สร้างประสบการณ์การเรียนรู้ที่มีทั้งความท้าทาย ความสมจริง และความปลอดภัยทางความรู้สึก
\end{enumerate}

% -----------------------------------------------------------
\section{วัตถุประสงค์}
% -----------------------------------------------------------

\begin{enumerate}[label=\arabic*)]
    \item เพื่อออกแบบและพัฒนาเกมเดินเรือด้วยระบบสั่งการด้วยเสียง (Voice-Controlled Ship Navigation Game) สำหรับใช้เป็นสื่อการเรียนรู้แบบดื่มด่ำ (Immersive Learning) ที่มีความสมจริงและช่วยลดความวิตกกังวลในการใช้ภาษาต่างประเทศของผู้เรียน
    \item เพื่อประเมินประสิทธิภาพของระบบและประเมินความพึงพอใจของผู้ใช้งาน โดยศึกษาความแม่นยำและความหน่วง (Latency) ของระบบปัญญาประดิษฐ์ในเกม รวมถึงประเมินผลสัมฤทธิ์ด้านประสบการณ์การใช้งานของผู้เล่นที่มีต่อการใช้เกมเป็นสื่อเพื่อการเรียนรู้ภาษา
\end{enumerate}

% -----------------------------------------------------------
\section{เป้าหมายและขอบเขต}
% -----------------------------------------------------------

\subsection{ขอบเขตด้านเนื้อหา}

ระบบที่พัฒนาเป็นเกมแนวจำลองสถานการณ์ (Simulation Game) แบบผู้เล่นคนเดียว (Single-player) ที่ใช้เสียงเป็นอุปกรณ์ควบคุมหลัก ครอบคลุมการทำงาน 5 ระบบหลัก ดังนี้

\begin{enumerate}[label=\arabic*)]
    \item \textbf{ระบบรู้จำเสียงและการประมวลผลคำสั่ง (Speech Recognition \& Command Parser System):} ระบบรับเสียงจากไมโครโฟน แปลงเป็นข้อความ และตรวจสอบความถูกต้องของคำสั่งภาษาอังกฤษ
    \item \textbf{ระบบฟิสิกส์และการควบคุมยานพาหนะ (Ship Navigation Physics):} ระบบจำลองการเคลื่อนที่ของเรือที่มีความเฉื่อยและการตอบสนองต่อพวงมาลัยที่สมจริง ทำให้ผู้เล่นต้องวางแผนล่วงหน้าในการออกคำสั่ง
    \item \textbf{ระบบลูกเรือปัญญาประดิษฐ์ (AI Crew System):} ระบบตัวละคร NPC ที่รับคำสั่งเสียงจากผู้เล่นและให้เสียงตอบกลับ เช่น การทวนคำสั่งเมื่อรับทราบ หรือการแจ้งเตือนเมื่อไม่เข้าใจคำสั่ง
    \item \textbf{ระบบสุ่มสร้างสิ่งกีดขวาง (Obstacle Generation System):} ระบบสร้างสิ่งกีดขวางแบบสุ่มในเส้นทางการเดินเรือ เพื่อท้าทายให้ผู้เล่นต้องตัดสินใจออกคำสั่งหลบหลีกในทันที
    \item \textbf{ระบบอินเทอร์เฟซผู้ใช้ (UI/HUD):} หน้าจอแสดงข้อมูลที่จำเป็น เช่น เรดาร์ตรวจจับสิ่งกีดขวาง เข็มทิศแสดงองศาเรือ มาตรวัดความเร็ว และสัญลักษณ์แสดงสถานะการรับเสียงของไมโครโฟน
\end{enumerate}

\subsection{ขอบเขตด้านประชากรและกลุ่มตัวอย่าง}

ทำการประเมินผลกับกลุ่มเป้าหมายคือนักศึกษาระดับปริญญาตรี หรือผู้ที่สนใจในการเรียนรู้และฝึกฝนทักษะการสื่อสารภาษาอังกฤษเบื้องต้นผ่านสื่อเกมแบบจำลองสถานการณ์ จำนวนกลุ่มผู้ทดลองประมาณ 40–50 คน

\subsection{ขอบเขตด้านเทคนิค}

\begin{enumerate}[label=\arabic*)]
    \item พัฒนาตัวเกมด้วย Game Engine Unity 3D และใช้ภาษา C\# สำหรับการพัฒนาสภาพแวดล้อมเกมและระบบฟิสิกส์ของการเดินเรือ
    \item นำเสนอภาพในรูปแบบมุมมองการเล่นแบบไอโซเมตริก (Isometric view) และรองรับการทำงานบนระบบปฏิบัติการ Windows PC
    \item ใช้ไมโครโฟนเป็นอุปกรณ์รับข้อมูล (Input) เสียงพูดภาษาอังกฤษ เป็นวิธีการหลักในการควบคุมทิศทางและความเร็วของเรือ โดยไม่มีการใช้คีย์บอร์ดหรือเมาส์ในการบังคับทิศทางโดยตรง
    \item ใช้ API หรือไลบรารีสำหรับการประมวลผลเสียง (Whisper) เพื่อแปลงเสียงพูดของผู้เล่นเป็นข้อความคำสั่ง (Text Parsing)
    \item ใช้บริการเชื่อมต่อระบบประมวลผล (API) สำหรับปัญญาประดิษฐ์ (OpenAI GPT-4o mini) เพื่อใช้ในระบบโต้ตอบและประมวลผลคำสั่งลูกเรือ
\end{enumerate}

% -----------------------------------------------------------
\section{ประโยชน์ที่คาดว่าจะได้รับ}
% -----------------------------------------------------------

\begin{enumerate}[label=\arabic*)]
    \item ผู้เล่นได้รับทักษะพื้นฐานด้านการสื่อสารภาษาอังกฤษและคำศัพท์ที่เกี่ยวข้องกับการเดินเรือ
    \item ผู้เล่นได้เกิดความมั่นใจในการพูดกล้าพูดภาษาอังกฤษมากยิ่งขึ้น
    \item สามารถนำไปใช้เป็นสื่อการสอนเพื่อส่งเสริมการเรียนรู้เชิงปฏิบัติ
\end{enumerate}